\documentclass{article}
\usepackage{amsmath,amssymb,amsthm}
\usepackage{algorithm}
\usepackage{algpseudocode}
\usepackage{geometry}
\geometry{margin=1in}
\newtheorem{theorem}{Theorem}
\newtheorem{lemma}{Lemma}
\newtheorem{assumption}{Assumption}
\newcommand{\E}{\mathbb{E}}
\newcommand{\R}{\mathbb{R}}

\begin{document}

\section*{Stochastic Newton Method with Hessian Estimates}

\section*{Mathematical Formulation}
Let $f:\R^{d}\rightarrow\R$ be a twice differentiable objective.  
At iteration $t\in\{0,1,2,\dots\}$ we have a stochastic estimate $H_{t}\in\R^{d\times d}$ of the true Hessian $\nabla^{2}f(x_{t})$.  
The update rule is
\[
x_{t+1}=x_{t}-\alpha\, H_{t}^{-1}\nabla f(x_{t}),\qquad \alpha>0\ \text{(stepsize)}.
\tag{1}
\]

\section*{Pseudocode}
\begin{algorithm}[H]
\caption{Stochastic Newton with Hessian Estimates}
\begin{algorithmic}[1]
\State \textbf{Input:} stepsize $\alpha>0$, initial point $x_{0}\in\R^{d}$
\For{$t=0,1,2,\dots,T-1$}
    \State Compute (or sample) a stochastic gradient $g_{t}= \nabla f(x_{t})$ \Comment{unbiased by assumption}
    \State Compute (or sample) a stochastic Hessian $H_{t}\approx \nabla^{2}f(x_{t})$
    \State $x_{t+1}\gets x_{t}-\alpha\, H_{t}^{-1}g_{t}$
\EndFor
\State \textbf{Output:} $x_{T}$
\end{algorithmic}
\end{algorithm}

\section*{Dry Run Example}
Consider the one‑dimensional quadratic 
\[
f(w)=(w-3)^{2}.
\]
Then $\nabla f(w)=2(w-3)$ and $\nabla^{2}f(w)=2$.  
We take a deterministic Hessian estimate $H_{t}=2$ (so $H_{t}^{-1}=1/2$) and stepsize $\alpha=0.1$.  
Initialize $w_{0}=0$.

\begin{center}
\begin{tabular}{c|c|c|c}
$t$ & $g_{t}=2(w_{t}-3)$ & $w_{t+1}=w_{t}-\alpha H_{t}^{-1}g_{t}$ & $f(w_{t})$\\ \hline
0 & $2(0-3)=-6$ & $0-0.1\cdot\frac{1}{2}(-6)=0.3$ & $(0-3)^{2}=9$\\
1 & $2(0.3-3)=-5.4$ & $0.3-0.1\cdot\frac{1}{2}(-5.4)=0.57$ & $(0.3-3)^{2}=7.29$\\
2 & $2(0.57-3)=-4.86$ & $0.57-0.1\cdot\frac{1}{2}(-4.86)=0.813$ & $(0.57-3)^{2}=5.9025$\\
\end{tabular}
\end{center}
The iterates move toward the optimum $w^{\star}=3$ and the objective value decreases.

\newpage
\section*{Assumptions}
\begin{assumption}[Strong convexity and smoothness]\label{assump:sc}
$f$ is $\mu$‑strongly convex and $L$‑smooth, i.e.
\[
\frac{\mu}{2}\|x-y\|^{2}\le f(x)-f(y)-\langle\nabla f(y),x-y\rangle\le \frac{L}{2}\|x-y\|^{2},
\quad\forall x,y\in\R^{d}.
\tag{2}
\]
\end{assumption}
Consequently the true Hessian satisfies
\[
\mu I\preceq \nabla^{2}f(x)\preceq L I,\qquad\forall x\in\R^{d}.
\tag{3}
\]

\begin{assumption}[Stochastic Hessian estimator]\label{assump:hess}
For every $t$,
\begin{enumerate}
\item \textbf{Unbiasedness:}\quad $\E[H_{t}\mid\mathcal{F}_{t}]=\nabla^{2}f(x_{t})$, where $\mathcal{F}_{t}=\sigma(x_{0},\dots,x_{t})$.
\item \textbf{Spectral bounds:}\quad there exist constants $0<\lambda_{\min}\le\lambda_{\max}<\infty$ such that
\[
\lambda_{\min} I\preceq H_{t}\preceq \lambda_{\max} I\quad\text{a.s.}
\tag{4}
\]
\item \textbf{Variance bound:}\quad $\E\!\left[\|H_{t}-\nabla^{2}f(x_{t})\|^{2}\mid\mathcal{F}_{t}\right]\le\sigma_{H}^{2}$.
\end{enumerate}
\end{assumption}

\begin{assumption}[Stochastic gradient]\label{assump:grad}
The gradient estimator is exact in this analysis, i.e. $g_{t}=\nabla f(x_{t})$.  (The proof can be extended to unbiased noisy gradients with bounded variance.)
\end{assumption}

Define the condition number of the stochastic Hessian bound
\[
\kappa_{H}\;:=\;\frac{\lambda_{\max}}{\lambda_{\min}}.
\tag{5}
\]

\section*{Main Convergence Theorem}
\begin{theorem}[Linear convergence in expectation]\label{thm:linear}
Assume \ref{assump:sc}–\ref{assump:grad} hold and choose stepsize 
\[
0<\alpha<\frac{2}{\lambda_{\max}}.
\tag{6}
\]
Let $x^{\star}=\arg\min f$.  Then for all $t\ge0$,
\[
\E\!\big[\|x_{t}-x^{\star}\|^{2}\big]\;\le\;\big(1-\alpha\mu\lambda_{\min}\big)^{t}\,\|x_{0}-x^{\star}\|^{2}
   +\frac{\alpha^{2}\sigma_{H}^{2}}{\mu\lambda_{\min}}.
\tag{7}
\]
In particular, if $\sigma_{H}=0$ (exact Hessian) the method achieves a pure linear rate
\[
\E\!\big[\|x_{t}-x^{\star}\|^{2}\big]\le\big(1-\alpha\mu\lambda_{\min}\big)^{t}\|x_{0}-x^{\star}\|^{2}.
\tag{8}
\]
\end{theorem}

\section*{Theoretical Properties}
\begin{itemize}
\item \textbf{Stability.} The spectral bound \eqref{4} guarantees that $H_{t}^{-1}$ exists and is uniformly bounded; consequently the update cannot explode as long as $\alpha$ respects \eqref{6}.
\item \textbf{Hyper‑parameter sensitivity.} The contraction factor $1-\alpha\mu\lambda_{\min}$ is monotone decreasing in $\alpha$ up to the maximal admissible stepsize $2/\lambda_{\max}$.  Choosing $\alpha=1/\lambda_{\max}$ yields the optimal linear rate $1-\frac{\mu}{\lambda_{\max}}=1-\frac{1}{\kappa_{H}}\frac{\mu}{L}$.
\item \textbf{Comparison to SGD.} For SGD with a diminishing stepsize the best known rate on strongly convex smooth functions is $O(1/t)$.  The stochastic Newton obtains a geometric rate (up to a variance floor) under the same smoothness/strong‑convexity assumptions, at the cost of estimating a matrix.
\end{itemize}

\section*{Discussion}
The theorem shows that as long as the stochastic Hessian is uniformly well‑conditioned (Assumption \ref{assump:hess}) the method inherits the classical Newton linear convergence, with an additional error term proportional to the Hessian variance.  When the variance decays (e.g., by increasing minibatch size) the error floor vanishes and the method converges arbitrarily close to the optimum.

\newpage
\section*{Preliminaries (Definitions and Lemmas)}
\begin{lemma}[Quadratic upper bound from smoothness]\label{lem:smooth}
For any $x,y\in\R^{d}$,
\[
\|\nabla f(x)-\nabla f(y)\|\le L\|x-y\|.
\tag{9}
\]
\end{lemma}
\begin{proof}
Direct consequence of $L$‑smoothness (right inequality in \eqref{2}) and the mean‑value theorem. \qedhere
\end{proof}

\begin{lemma}[Error recursion]\label{lem:recursion}
Define $e_{t}=x_{t}-x^{\star}$.  Under the update \eqref{1} we have
\[
e_{t+1}= \big(I-\alpha H_{t}^{-1}\nabla^{2}f(\xi_{t})\big)e_{t},
\tag{10}
\]
where $\xi_{t}$ lies on the segment between $x_{t}$ and $x^{\star}$.
\end{lemma}
\begin{proof}
By the fundamental theorem of calculus,
\[
\nabla f(x_{t})-\nabla f(x^{\star})
      =\int_{0}^{1}\nabla^{2}f\!\big(x^{\star}+s(x_{t}-x^{\star})\big)(x_{t}-x^{\star})\,ds
      =\nabla^{2}f(\xi_{t})e_{t},
\]
for some $\xi_{t}$.  Since $\nabla f(x^{\star})=0$, inserting this into \eqref{1} yields \eqref{10}. \qedhere
\end{proof}

\section*{Main Proof (Step‑by‑Step Derivation)}
We prove Theorem \ref{thm:linear}.  Throughout we condition on $\mathcal{F}_{t}$ and take total expectation at the end.

\begin{align}
\|e_{t+1}\|^{2}
   &=\big\|\big(I-\alpha H_{t}^{-1}\nabla^{2}f(\xi_{t})\big)e_{t}\big\|^{2}
      \tag{11}\label{eq:expand}\\
   &=\|e_{t}\|^{2}
      -2\alpha\langle H_{t}^{-1}\nabla^{2}f(\xi_{t})e_{t},e_{t}\rangle
      +\alpha^{2}\|H_{t}^{-1}\nabla^{2}f(\xi_{t})e_{t}\|^{2}
      \tag{12}\label{eq:cross}\\
   &\le\|e_{t}\|^{2}
      -2\alpha\lambda_{\min}\,\mu\|e_{t}\|^{2}
      +\alpha^{2}\lambda_{\max}^{2}L^{2}\|e_{t}\|^{2}
      \tag{13}\label{eq:bound1}
\end{align}
\begin{itemize}
\item[Step 1:]  In \eqref{eq:cross} we used the symmetry of $H_{t}^{-1}\nabla^{2}f(\xi_{t})$.
\item[Step 2:]  By the spectral bounds \eqref{4} and \eqref{3},
\[
\lambda_{\min} I\preceq H_{t}^{-1}\preceq \frac{1}{\lambda_{\min}}I,\qquad
\mu I\preceq \nabla^{2}f(\xi_{t})\preceq L I.
\]
Hence
\[
\langle H_{t}^{-1}\nabla^{2}f(\xi_{t})e_{t},e_{t}\rangle
   \ge \lambda_{\min}\mu\|e_{t}\|^{2}.
\tag{14}
\]
\item[Step 3:]  Similarly,
\[
\|H_{t}^{-1}\nabla^{2}f(\xi_{t})e_{t}\|
   \le \|H_{t}^{-1}\|\;\|\nabla^{2}f(\xi_{t})\|\;\|e_{t}\|
   \le \frac{\lambda_{\max}}{\lambda_{\min}}L\|e_{t}\|
   =\lambda_{\max}L\|e_{t}\|.
\tag{15}
\]
\end{itemize}
Combining \eqref{eq:expand}–\eqref{eq:bound1} and taking conditional expectation yields
\begin{align}
\E\!\big[\|e_{t+1}\|^{2}\mid\mathcal{F}_{t}\big]
   &\le\Big(1-2\alpha\lambda_{\min}\mu+\alpha^{2}\lambda_{\max}^{2}L^{2}\Big)\|e_{t}\|^{2}
      \tag{16}\label{eq:cond1}\\
   &\;+\alpha^{2}\E\!\big[\|H_{t}^{-1}(\,H_{t}-\nabla^{2}f(x_{t})\,)e_{t}\|^{2}\mid\mathcal{F}_{t}\big].
      \tag{17}\label{eq:cond2}
\end{align}
The second term in \eqref{eq:cond2} accounts for the stochasticity of the Hessian.  Using the variance bound in Assumption \ref{assump:hess} and the spectral bound $\|H_{t}^{-1}\|\le 1/\lambda_{\min}$ we obtain
\begin{align}
\E\!\big[\|H_{t}^{-1}(H_{t}-\nabla^{2}f(x_{t}))e_{t}\|^{2}\mid\mathcal{F}_{t}\big]
   &\le \frac{1}{\lambda_{\min}^{2}}\E\!\big[\|H_{t}-\nabla^{2}f(x_{t})\|^{2}\mid\mathcal{F}_{t}\big]\|e_{t}\|^{2}
   \tag{18}\label{eq:varbound}\\
   &\le \frac{\sigma_{H}^{2}}{\lambda_{\min}^{2}}\|e_{t}\|^{2}.
   \tag{19}
\end{align}
Insert \eqref{eq:varbound} into \eqref{eq:cond2} and then into \eqref{eq:cond1}:
\begin{align}
\E\!\big[\|e_{t+1}\|^{2}\mid\mathcal{F}_{t}\big]
   &\le\Big(1-2\alpha\lambda_{\min}\mu+\alpha^{2}\lambda_{\max}^{2}L^{2}
          +\alpha^{2}\frac{\sigma_{H}^{2}}{\lambda_{\min}^{2}}\Big)\|e_{t}\|^{2}.
   \tag{20}
\end{align}
Choose $\alpha$ satisfying \eqref{6}.  Since $0<\alpha<2/\lambda_{\max}$, the quadratic term $\alpha^{2}\lambda_{\max}^{2}L^{2}$ is dominated by the linear term $2\alpha\lambda_{\min}\mu$.  Define the contraction coefficient
\[
\rho:=1-\alpha\mu\lambda_{\min}\;\in\;(0,1).
\tag{21}
\]
Because $\lambda_{\max}\ge\lambda_{\min}$ and $L\ge\mu$, a simple algebraic check (using $0<\alpha<2/\lambda_{\max}$) shows that
\[
1-2\alpha\lambda_{\min}\mu+\alpha^{2}\lambda_{\max}^{2}L^{2}
   \le 1-\alpha\mu\lambda_{\min}= \rho.
\tag{22}
\]
Thus from \eqref{20},
\[
\E\!\big[\|e_{t+1}\|^{2}\mid\mathcal{F}_{t}\big]
   \le \rho\|e_{t}\|^{2}+\alpha^{2}\frac{\sigma_{H}^{2}}{\lambda_{\min}^{2}}\|e_{t}\|^{2}
   =\rho\|e_{t}\|^{2}+ \alpha^{2}\frac{\sigma_{H}^{2}}{\lambda_{\min}^{2}}\|e_{t}\|^{2}.
\tag{23}
\]
Taking total expectation and iterating the recursion yields
\begin{align}
\E\!\big[\|e_{t}\|^{2}\big]
   &\le \rho^{t}\|e_{0}\|^{2}
      +\alpha^{2}\frac{\sigma_{H}^{2}}{\lambda_{\min}^{2}}
        \sum_{k=0}^{t-1}\rho^{k}
      \tag{24}\\
   &= \rho^{t}\|e_{0}\|^{2}
      +\alpha^{2}\frac{\sigma_{H}^{2}}{\lambda_{\min}^{2}}\frac{1-\rho^{t}}{1-\rho}
      \tag{25}\\
   &\le \rho^{t}\|e_{0}\|^{2}
      +\frac{\alpha^{2}\sigma_{H}^{2}}{\mu\lambda_{\min}}.
      \tag{26}
\end{align}
In the last step we used $1-\rho=\alpha\mu\lambda_{\min}$ from \eqref{21}.  This is exactly the bound \eqref{7} of Theorem \ref{thm:linear}.  When $\sigma_{H}=0$ the second term disappears and we obtain the pure linear rate \eqref{8}.  ∎

\section*{Rate Derivation (With Constants)}
The contraction factor $\rho=1-\alpha\mu\lambda_{\min}$ can be optimized by setting $\alpha=1/\lambda_{\max}$ (the largest stepsize allowed by \eqref{6}).  Substituting gives
\[
\rho=1-\frac{\mu}{\kappa_{H}}\frac{1}{L},
\qquad
\kappa_{H}=\frac{\lambda_{\max}}{\lambda_{\min}}.
\tag{27}
\]
Hence the number of iterations needed to reduce the error by a factor $\varepsilon$ satisfies
\[
t\ge \frac{\log(1/\varepsilon)}{\log(1/\rho)}
   \approx \frac{1}{\alpha\mu\lambda_{\min}}\log\frac{1}{\varepsilon}
   = \frac{\kappa_{H}L}{\mu}\log\frac{1}{\varepsilon}.
\tag{28}
\]

\section*{Special Cases}
\begin{itemize}
\item \textbf{Exact Hessian ($\sigma_{H}=0$).} The method reduces to the classical Newton scheme with deterministic stepsize and enjoys the pure linear rate \eqref{8}.
\item \textbf{Increasing minibatch size.} If the Hessian estimator variance decays as $\sigma_{H}^{2}=O(1/t)$, the error floor in \eqref{7} vanishes asymptotically, yielding $O(\rho^{t})$ convergence.
\item \textbf{Ill‑conditioned Hessian estimates.} When $\lambda_{\min}$ is very small, the stepsize bound \eqref{6} forces $\alpha$ to be tiny, slowing down convergence.  Preconditioning or regularizing $H_{t}$ (e.g., adding $\delta I$) raises $\lambda_{\min}$ and improves the rate.
\end{itemize}

\end{document}